\documentclass[../main.tex]{subfiles}
\begin{document}
%==========================================TIÊU ĐỀ TẬP========================================
\begin{table}[h]
    \begin{adjustwidth}{-1cm}{}
        \begin{tabular}{>{\centering\arraybackslash}p{6cm}|>{\raggedright\arraybackslash}p{12.5cm}}
            \multirow{5}{6cm}
            % Cột bên trái: ÉPISODE và số tập
            {\\[-40pt]
                \flaregothic\fontsize{20pt}{20pt}\selectfont \centering
                \color{eptype}HISTOIRE\color{black}\\
                \freshbot\fontsize{72pt}{72pt}\selectfont
                \color{epnum}15\\[6pt]
                \cafeteria\fontsize{20pt}{20pt}\selectfont\color{epnum}(D-15)\color{black}
                \\[18pt]
                % \flaregothic\fontsize{14pt}{14pt}\selectfont
                % \color{epattr}BIENVENUE, \uline{\color{Banpire} Mel} !
                \sen\fontsize{18pt}{18pt}\selectfont
                \color{epattr}Mini-histoire
                \color{black}
            }
             &
            % Dòng thứ nhất: tên tập tiếng Nhật
            {\fotskip\fontsize{18pt}{18pt}\selectfont ここは、いつでも\ruby{扉}{とびら}が\ruby{開}{あ}いている}
            \\[6pt]
            % Dòng thứ hai: tên tập tiếng Anh
             & {\cafeteria\fontsize{18pt}{18pt}\selectfont Because Here, the Door is Always Open} \\[4pt]
            % Dòng thứ ba: tên tập tiếng Việt
             & {\vnmsans\fontsize{18pt}{18pt}\selectfont Những ngày đầu của cặp chị em nhà Mèo}                      \\[8pt]
            % òng thứ tư: Nhân vật xuất hiện
             & {\caviar{\fontsize{14pt}{14pt}\selectfont Nhân vật: Theo từng câu chuyện}}
            \\[8pt]
            % Dòng cuối cùng: Release Date
            % & {\continuum\fontsize{16pt}{16pt}\selectfont Release Date: 07/02/2021}\\[18pt]
        \end{tabular}
    \end{adjustwidth}
\end{table}
%=========================================TRUYỆN CHÍNH========================================
\color{black}

\txtbx[2]{
    {\color{vol} \uline{Tóm tắt:}} Sáu mẩu chuyện nhỏ về những ngày đầu của Mikeneko và Koneko tại ký túc xá Hikawa: từ việc làm quen với bản đồ và quy tắc nhà, cuộc thi nấu cơm cuộn bị Suzuno "xử đẹp", màn cày game đêm đầy toxic với Ryuuhou, khoảnh khắc ấm áp trên ban công với Naomi và kẹo mút, sự chu đáo của Suzuno khi dọn bếp, và lời tâm sự của Mikeneko về nỗi sợ bị từ chối, rồi Kuon trấn an rằng cánh cửa luôn mở cho hai chị em.
}

\stpt{-Histoire 1-}
\stcap{Chiếc bản đồ và quy tắc}
\begin{center}
    {\caviar\fontsize{12pt}{12pt}\selectfont Nhân vật: \emp{kaigou2} \emp{mike1} \emp{kson}}
\end{center}

Nắng sớm của mùa đông len lỏi qua ô cửa kính ở cuối hành lang tầng năm, vắt những vệt sáng vàng nhạt lên sàn nhà gỗ. Không khí bình minh còn vương chút lạnh giá, chỉ có tiếng bước chân lẹp kẹp vang lên đều đặn, xen lẫn chút ngần ngừ.

Mikeneko đứng đóng băng trước ngã rẽ hành lang, hai tay siết chặt một tờ giấy A4 vốn phẳng phiu nay đã nhăn nhúm như bị nhào nặn hàng chục lần. Trên mặt giấy là những đường bút bi xanh nguệch ngoạc, bản đồ ký túc xá và khu vực xung quanh mà cô nàng tự phác thảo sau chuyến xe tối qua, đầy dấu chấm hỏi và mũi tên lạc lối.

Dù cô chỉ cần mở miệng nhờ tìm cửa hàng tiện lợi để mua đồ vệ sinh cá nhân và vài lon nước ngọt, cái tự tôn cao ngút trời của cô nàng Mèo hồng lại không cho phép. Mới chuyển đến tầng tám hôm qua, nếu giờ lại lẽo đẽo xuống hỏi người dưới đường đi ra cửa hàng tạp hóa, khác gì cô thừa nhận mình hoàn toàn lạc lối? Mà Mikeneko của người ta, làm sao có thể thừa nhận điểm yếu chứ?

Cô đứng ngẩn ngơ hàng mươi phút, lông mày nhíu chặt thành một đường, khuôn mặt xoay vòng giữa bực bội và ngượng ngùng. Tờ giấy bị cô xoay 180 độ rồi lại xoay ngang, môi lầm bầm những câu rủa sả nho nhỏ vào cái bản đồ vô dụng. ``\textit{Làm sao mà ngã tư này lại trông giống mọi ngã tư khác thế nhỉ? Rốt cuộc lối nào mới là lối ra đây?}''

``Này, Mike-paisen, chị đứng làm quân kỳ ở đây từ bao giờ thế?''

Một giọng nói trêu chọc vang lên từ phía sau, phá vỡ sự im lặng của hành lang. Kson bước ra từ căn hộ tầng năm của mình, một tay cầm cốc cà phê đen nghi ngút khói, mái tóc còn rũ rượi vì vừa mới ngủ dậy. Nhìn thấy vẻ mặt ``thua trận'' căng như dây đàn của người tiền bối, nàng Tổng trưởng không kìm được nụ cười tủm tỉm hiện ra. Thật ra cô đã đứng đây quan sát được mười phút rồi. ``Chắc chị lại lạc đường nữa nhỉ?''

Mikeneko giật thót người, nhanh như chớp giấu tờ giấy sau lưng, lưng đập thẳng tắp như cột điện. Tai mèo trên đầu nàng xù lên một chút vì bất ngờ. ``A-Ai làm quân kỳ chứ?! Chị chỉ\dots\ đang thẩm định lại sơ đồ không gian và lối thoát hiểm của tòa nhà thôi! Đừng có mà suy đoán lung tung!''

Kson nhấp một ngụm cà phê từ từ, ánh mắt đùa giỡn nheo lại như mèo vờn chuột. ``Thẩm định sơ đồ mà vẽ cái ngã tư to đùng lại đâm thẳng vào tường nhà kho thế kia à? Cần chỉ đường ra ngoài mua đồ thì nói một tiếng, việc gì phải chịu khổ thế này, hả chị?''

``Chị không cần! Chị biết thừa đường đi!'' nàng Mèo hồng phồng má, giọng sặc mùi chống chế nhưng hai má đã đỏ bừng, lan lên tận vành tai mỏng. Tay nàng siết chặt tờ giấy sau lưng, gần như muốn xé nát nó để che giấu sự bối rối.

``Biết đường mà đứng quay mòng mòng ở đây vậy từ hồi sáng?''

Tiếng bước chân nặng thịch vang lên từ hướng cầu thang bộ, cắt ngang cuộc đối đầu nho nhỏ của hai người. Kaigou bước lên, trên vai vác một bao cát tập thể dục nặng năm mươi cân như chỉ là chiếc gối bông nhẹ tênh, dù tóc cô vẫn đang buộc gọn gàng sau khi tập thể dục. Giọt mồ hôi lăn nhẹ trên gò má sắc sảo của cô chị Hai nhà Hikawa, thêm vào vẻ ngoài nghiêm nghị thường ngày. ``Hai đứa đứng đây làm gì mà sớm thế?''

Nàng Tổng dừng lại, ánh mắt lạnh như tiền lướt qua Mikeneko rồi quay sang Kson. Cô tặc lưỡi một cái rõ to, hạ bao cát xuống sàn tạo ra tiếng *RẦM* nhe, khiến Mikeneko giật mình thót tim, tai mèo xù lên hết cỡ.

``Chị lại còn tự vẽ bản đồ cơ đấy,'' Kaigou khoanh tay, giọng điệu thản nhiên không chút nịnh nọt. ``Khu này ngõ ngách chằng chịt, cầm cái tờ giấy nguệch ngoạc đó đi ra ngoài có mà mất tích đến tối. Nếu chị định mua đồ thì ít nhất cũng phải lên tiếng để người ta còn chỉ chứ. Đứng đây cả sáng cũng không ra được đường đâu.''

``Đó! Kaigou-paisen cũng nói thế rồi còn gì!'' Kson lập tức hưởng ứng lời của người chị Hai, giọng đầy vẻ ``ta cùng phe''. Cô nhấp thêm một ngụm cà phê, rướn người lên chờ xem màn trình diễn của cựu tiền bối.

Mikeneko mấp máy môi định cãi lại, nhưng Kaigou đã điềm tĩnh ngắt lời bằng phong thái thực dụng vốn có. ``Quanh đây chỉ có hai chỗ để mua đồ. Nếu muốn mua mấy thứ lặt vặt như kim chỉ, băng cá nhân hay đồ ăn vặt truyền thống thì ra khỏi cửa chính, rẽ phải đi hết đầu phố. Ở đó có tiệm tạp hóa của một bác cô, giá rẻ nhưng tầm 11 giờ đêm là đóng cửa rồi, cũng không có mấy đồ có sẵn đâu.''

Nàng Tổng ngoắc tay ra hướng ngược lại, chỉ rõ ràng từng bước đi. ``Còn nếu rẽ trái ra phía ngã tư cách đây tầm bốn trăm mét, có một cửa hàng tiện lợi 24/7. Ở đó có máy pha cà phê tự động, đồ ăn đóng hộp, bàn ghế ngồi lại và đầy đủ thiết bị cá nhân. Muốn nhanh gọn thì đi hướng ngã tư, đỡ mất thời gian.''

Nàng Mèo hồng chớp chớp mắt, hai tai mèo trên đầu khẽ giật giật như thể nhận được tín hiệu. Cô cuống quýt rút tờ giấy sau lưng ra, lén lút ghi chép lại những thông tin đó bằng tốc độ ánh sáng, đầu bút bi xiên xiên trên mặt giấy.

``Mà, Kaigou-paisen này, chị đang làm gì đấy ạ?'' Kson hỏi, tay cầm cốc cà phê đã vơi đi gần một nửa.

Kaigou xóc lại bao cát trên vai, quay người bước đi về phía cầu thang: ``Chị chuẩn bị xuống sảnh dọn dẹp và kiểm tra lại hệ thống nước nóng. Mikeneko-san nếu muốn thì đi cùng luôn, chị cùng có mấy điều muốn dặn dò đây.''

Mikeneko ngập ngừng một chút, nhưng sự thực tế đã chiến thắng cái tự tôn vô lý trong người. Cô lủi thủi bước theo sau nàng Tổng, giữ một khoảng cách chừng hai bước chân, như học sinh đi sau giáo viên. Kson thì thong thả đi sóng đôi bên cạnh, vừa đi vừa tiếp tục nhâm nhi cốc cà phê chưa hết.

``Quy củ ở nhà Hikawa này cũng nghiêm lắm đấy,'' giọng Kaigou vang lên đều đặn, chuyên nghiệp và rõ ràng như một quản lý đang hướng dẫn nhân viên mới vào nghề. ``Buổi sáng, nếu muốn tắm rửa thì phải làm sau tầm 7 giờ 30 sáng. Khoảng giờ trước đó là thời điểm bọn trẻ - Ryuu, Suzu và Nao - chuẩn bị đi học và Aniki chuẩn bị đi làm, bình nước nóng tầng trung tâm sẽ bị quá tải. Tắm tối thì tự do, nhưng không khuyến khích tắm đêm sau 11 giờ, gây ảnh hưởng người khác.''

Nàng Mèo hồng vừa bước vừa cắm cúi lấy đầu bút bi ghi chép lia lịa lên mép tờ giấy A4, mắt không rời khỏi mặt giấy như sợ bỏ sót một chữ nào. ``\textit{Tắm sáng sau 7h30\dots\ Mình nhớ được rồi\dots}''

``Về phần bếp chung ở tầng sáu,'' Kaigou tiếp tục liệt kê những quy định quan trọng, giọng nói không hề ngắt quãng, ``đồ ăn mua về tủ lạnh phải dán nhãn tên mình. Món nào không dán nhãn sau ba ngày sẽ bị Suzu dùng để làm nguyên liệu nấu ăn, nên trừ khi đồ mà mấy đứa mua về xác định dùng chung cho cả nhà thì phải nhớ dán nhãn. Giặt đồ thì phòng cô có máy giặt riêng trên tầng sáu này, nhưng tuyệt đối không phơi đồ ra ban công phía trước vào thứ Ba và thứ Năm vì đó là ngày xe thu gom rác của khu vực đi qua, bụi lắm.''

Kson đi bên cạnh khẽ cười hơ hơ, chen vào bằng chất giọng bộc bạch, thêm vào những lưu ý của Kaigou. ``Tóm lại là có hai quy tắc vàng sinh tồn ở cái nhà này mà cô phải thuộc lòng, Mikeneko-paisen à. Quên một trong hai là chịu trận đấy.''

Mikeneko ngước lên, ánh mắt đầy tò mò bỏ khỏi tờ giấy, lập tức hỏi liền. ``Quy tắc gì? Chị nói nhanh để em ghi lại.'' Cô nàng cầm sẵn bút bi, chuẩn bị viết liền những điều quan trọng này.

Nàng Tổng trưởng đưa một ngón tay gun lên chỉ trỏ, nhấn mạnh từng quy tắc như sợ Mikeneko quên. ``Thứ nhất: Đừng bao giờ chạm vào ngăn tủ lạnh có dán băng keo đen ghi tên Ryuu-chan. Nó mà thấy mất một lon nước tăng lực hay gói bim bim nào trong đó là nó biến thành `quỷ toxic' sấy chị cháy máy đấy. Thứ hai: Mấy món bánh ngọt hay tráng miệng có đánh dấu nốt ruồi hình trái tim của Suzu ấy, đừng dại mà ăn trộm. Em ấy sẽ không chửi đâu, nhưng sẽ nhìn em bằng đôi mắt dị sắc rấn rấn nước mắt, và tin chị đi, lúc đó đến cả Aniki cũng chẳng cứu nổi em khỏi sự trừng phạt của lương tâm đâu.''

Nàng Mèo hồng rùng mình một cái, thầm ghi thêm hai dòng cảnh báo đỏ rực vào góc sổ, bôi đậm cả hai quy tắc kia để không bao giờ quên. ``\textit{Trời ơi, nhà này có nhiều quy định cấm kỵ thế sao?}'' 

Đi xuống đến sảnh tầng trệt, Kaigou dừng lại chỉ tay ra phía cánh cửa kính lớn, chỉ rõ hai lối đi một lần nữa. ``Đấy, lối ra ở kia. Rẽ phải là tiệm tạp hóa đầu phố, rẽ trái là ngã tư cửa hàng tiện lợi. Chị cần gì thì tự đi mà mua. Không thì nhờ Kson-san trở đi ấy. Chị cũng phải chuẩn bị đi làm bây giờ đây.''

Nói rồi, nàng Tổng vác bao cát đi thẳng ra sân sau mà không chờ một lời cảm ơn, đúng phong thái coi việc hỗ trợ là nghĩa vụ hiển nhiên chứ chẳng phải ban phát ân huệ. Kaigou luôn như thế, lặng lẽ giúp đỡ mọi người trong nhà mà không bao giờ đòi hỏi đáp trả. 

``Mấy đứa nhớ đi nhanh về sớm nhé,'' cô vẫy tay nhẹ trước khi biến mất.

Mikeneko đứng ở sảnh lớn, nhìn tờ giấy A4 trong tay mình. Từ một bản đồ nguệch ngoạc sai bét, giờ đây nó đã được phủ kín bởi những dòng ghi chú tỉ mỉ: giờ tắm nóng lạnh, lịch thu gom rác, hai hướng đi mua đồ, và cả những ``vùng cấm địa'' trong tủ lạnh. Tất cả mọi thứ rõ ràng hơn bao giờ hết.

Cô nàng lầm bầm điều gì đó trong miệng, hai tai mèo khẽ cọ vào nhau đầy bối rối. Nơi này\dots\ đông đúc, rắc rối, ồn ào và đầy những con người kỳ lạ. Thế nhưng, đằng sau sự hỗn loạn ấy lại là một bộ máy vận hành cực kỳ kỷ luật, ngăn nắp và có quy củ riêng. Không giống bất kỳ nơi nào cô từng sống trước đây. Nó ấm áp hơn rất nhiều.

Nó khác hoàn toàn với căn hộ chung cư lạnh lẽo, tĩnh mịch mà cô từng sống một mình trong suốt những năm tháng cô đơn trước đây. Không ai chỉ đường cho cô, không ai bắt cô tuân theo giờ giấc, cũng chẳng ai cảnh báo cô về những chiếc tủ lạnh bị rút phích hay mì gói hết đát. Cuộc sống trước đây của cô chỉ có một mình cô, không một ai quan tâm. Giờ thì mọi thứ đã khác.

Kson cười mỉm, vỗ nhẹ vào vai Mikeneko một cái rồi rảo bước đi trước về phía cửa ra vào: ``Cần dẫn đường ra ngã tư không, tiền bối? Em cũng định ra đó mua đồ đây. Đi cùng cũng vui mà.''

``C-Chị tự đi được!'' nàng Mèo hồng cau mày gắt nhẹ, nhưng đôi tai mèo trên đầu lại khẽ vẫy vẫy đầy vui vẻ, nhanh chân bước theo sau cậu. Mặt nàng vẫn còn hơi đỏ, nhưng trong lòng đã thấy ấm áp lạ thường. ``Chị không cần ai dẫn đường đâu!''

Mikeneko thở ra một hơi dài, cẩn thận gấp tờ giấy lại cho vào túi áo, rồi cất bước đi về phía cửa hàng tiện lợi ở ngã tư. Lần này, bước chân của cô nàng đã hoàn toàn tự tin và không còn ngập ngừng nữa. Ánh nắng sớm chiếu lên vai nàng, mang theo chút ấm áp của ngày mới, và của một cuộc sống mới đầy bất ngờ ở nhà Hikawa.

\newpage

\stpt{-Histoire 2-}
\stcap{Cuộc thi nấu ăn}
\begin{center}
    {\caviar\fontsize{12pt}{12pt}\selectfont Nhân vật: \emp{suzuno2} \emp{mike1} \emp{delu}}
\end{center}

Buổi chiều muộn, ánh hoàng hôn nhuộm cả khung cửa sổ bếp tầng sáu thành một màu cam ấm áp, ánh nắng còn sót lại len lỏi rải trên mặt bàn đá. Không khí trong căn bếp chật chội này bất chợt trở nên căng thẳng không khác gì sợi dây đàn đang bị kéo căng đến mức tưởng chừng đứt.

Mikeneko đứng thẳng trước bếp, hai tay cầm chồng lá rong biển khô và cái tô cơm trắng vừa xới xong, hơi nóng còn bay nghi ngút lên má. Cô đã nhẩm tính cả tuần qua để có dịp thể hiện, và hôm nay chính là ngày cô chuẩn bị làm một bữa cơm ngon lành để ``trả ơn'' Kuon, vì dù sao đi nữa, cứ ở nhờ nhà người ta mà để anh ấy lo hết mọi việc thì cái lòng tự trọng cao ngất trời của cô không cho phép.

Nhưng ngay khi cô vừa chuẩn bị dồn toàn bộ tâm sức vào món ăn, một bóng người khác đứng hiện ra chỉ cách cô chưa đầy một mét. Delutaya, cầm trên tay đúng\dots\ chồng lá rong biển y hệt, và cả một tô cơm trắng cũng còn đang bốc khói nóng hổi.

Cô nàng Tam giác cũng chẳng muốn bị ai xem là kẻ chỉ biết ăn sẵn. Quan trọng hơn, vết thương lòng từ lần đầu tiên nấu ăn cho cả nhà vẫn còn đau: Kuon đã nếm thử món cô làm rồi thản nhiên nhận xét, ``So với đồ bố anh nấu thì vẫn còn kém xa.'' Lần này cô phải lấy lại sĩ diện, bằng mọi giá.

Hai người bật ngửa nhìn nhau. Cả hai cùng cúi xuống kiểm tra bàn bếp, và nhận ra nguyên liệu của đối phương giống hệt nhau từng li từng tí: dưa chuột thái sợi mảnh, trứng rán cắt thanh dài, miếng xá xíu thơm lừng, cà rốt bào sợi và cả thanh cua dai ngon.

Sự im lặng bao trùm cả căn bếp trong đúng năm giây dài như cả thế kỷ. Không ai chịu mở miệng trước, cũng chẳng ai có ý định nhường lại mặt bếp cho người kia.

``Cô\dots\ cũng định làm cơm cuộn à?'' Mikeneko nhướn mày cao lên, khoanh tay trước ngực như đang tra hỏi kẻ đột nhập.

``Thì sao? Bếp này là của chung mà, cô có ý kiến gì hả?'' Delutaya cũng hất cẳng đáp trả, mắt chẳng dời khỏi khay nguyên liệu của mình thêm một giây nào.

Không cần lời tuyên bố, không cần báo hiệu. Một cuộc chiến ngầm chính thức bùng nổ giữa hai người. Cả hai cùng ngầm hiểu rằng sẽ thi nhau nấu cùng lúc, ai làm ra đĩa cơm cuộn ngon hơn sẽ chính là ``đầu bếp số một'' của gian bếp hôm nay.

Đúng vào lúc không khí đang căng nhất, cánh cửa bếp lại bật mở. Suzuno đeo chiếc tạp dề in hình hoa nhỏ bước vào, tay ôm cái rổ rau củ vừa mua từ siêu thị về, định chuẩn bị bữa tối cho cả nhà. Vừa đặt chân qua ngưỡng cửa, cô chớp chớp cặp mắt dị sắc đỏ-xanh ngơ ngác, chẳng hiểu sao hai chị lại đứng sừng sỏ đối mặt nhau như sắp sửa cãi nhau.

``Ơ\dots\ Mikeneko-san? Delutaya-san? Hai chị đang làm gì ở đây vậy?''

``Suzuno-chan! Đúng lúc quá!'' Mikeneko liền vớ lấy cơ hội, kéo cô bé lại gần cái bàn bếp, giọng đầy phấn khích. ``Em làm trọng tài cho bọn chị nhé! Xem ai cuộn cơm ngon hơn!''

``Hả? Nhưng em\dots'' Suzuno vẫn còn đang ngẩn ngơ chưa kịp hiểu đầu cua tai nheo ra sao, thì Delutaya đã nhanh chóng ấn một chiếc thìa sạch và một đôi đũa vào tay em.

``Nhờ em đấy, Suzuno-chan. Cứ chấm điểm thật công tâm vào, đừng có thiên vị ai đâu nhé.''

Thế là cô chị Ba nhà Hikawa vô tình bị đẩy vào ghế ``trọng tài tối cao'' mà chẳng có cơ hội từ chối, đầu óc vẫn còn quay mòng mòng chẳng hiểu chuyện gì đang xảy ra với căn bếp nhà mình.

Chẳng mấy chốc, căn bếp nhỏ biến thành một chiến trường lộn xộn thực thụ.

Mikeneko thuộc kiểu người hành động nhanh như chớp, tay rải cơm trên lá rong biển không ngơi nghỉ, nhưng vì tính khí hấp tấp nên cứ vài phút lại quên mất công thức, lại phải ngoái đầu gọi Suzuno.

``Suzuno-chan! Cho giấm gạo vào cơm tỉ lệ bao nhiêu là đúng ấy nhỉ?!''

``Chết rồi, lượng muối mình nếm thử có vẻ hơi mặn, có bị quá không em?!''

``Rong biển thì nên đặt mặt nhám hay mặt bóng lên trên cơ chứ?!''

Còn Delutaya thì lại đi theo trường phái tỉ mỉ đến từng milimét. Cô đong từng thìa gia vị một cách chính xác, xếp từng sợi cà rốt trên mặt cơm thẳng tắp như được đóng khung. Nhưng yếu điểm của cô là\dots\ cứ vài giây lại phải lén liếc sang bên kia, xem Mikeneko đã làm đến đâu rồi, sợ rằng đối phương sẽ vượt mặt mình.

``\textit{Hừm, cô ta lại rắc thêm nhiều vừng đen thế à? Hay là mình cũng cho thêm một ít cho ngon?}'' cô vừa cắn môi suy nghĩ, vừa lén lút lấy lọ vừng trên kệ rắc thêm vào phần cơm của mình.

Suzuno bị kẹp ở giữa, chạy đôn chạy đáo hết bên này sang bên kia. Bản tính hiền lành khiến cô không muốn làm ai buồn, nên cứ hễ Mikeneko gọi là cô lật đật chạy qua giúp nêm gia vị, rồi Delutaya hỏi gì lại vội vàng bước qua chỉnh giúp. Cô nàng cứ như con thoi nhỏ bay qua lại giữa hai bếp lửa.

Sau hơn nửa tiếng ``chiến đấu'' căng thẳng, hai đĩa cơm cuộn hoàn chỉnh cuối cùng cũng được đặt lên bàn.

Mikeneko tự hào vỗ ngực, đĩa cơm của cô trông vô cùng hấp dẫn, nhân đầy ắp trào ra cả bên ngoài. Còn Delutaya thì nhẹ nhàng vuốt lại vết nhăn trên tạp dề, đĩa cơm của cô được cắt khoanh đều tăm tắp, xếp thành hình hoa cúc đẹp mắt chẳng khác gì hàng ngoài tiệm.

``Mời trọng tài thưởng thức ạ!'' cả hai cùng đồng thanh nói, mắt rạng ngời chờ đợi.

Nói về khoản ăn uống, Suzuno chính là chuyên gia hàng đầu không ai tranh được của nhà Hikawa. Vẻ ngây ngô thường ngày lập tức biến mất, gương mặt cô nàng trở nên nghiêm túc lạ thường. Cô cầm đôi đũa, cẩn thận gắp một miếng cơm của Mikeneko đưa vào miệng, nhai chậm rãi, sau đó lại nếm tiếp một miếng của Delutaya.

Hai người đứng phía sau nín thở, tim đập mạnh như sắp nhảy ra ngoài, hồi hộp chờ đợi lời phán quyết cuối cùng.

Suzuno nuốt miếng cuối cùng, chắp hai tay trước ngực rồi trầm ngâm đưa ra nhận xét đầy chuyên môn: ``Cả hai\dots\ đều ngon theo cách riêng của mình ạ. Cơm cuộn của Mikeneko-san nêm nếm đậm đà, ăn rất đã miệng, nhưng tay cuộn hơi lỏng nên nhặt lên dễ bị rơi nhân ra. Còn của Delutaya-san thì cuộn chắc chắn, trình bày đẹp mắt, vị vừa ăn, chỉ là hơi thiếu một chút muối để làm nổi bật vị ngọt tự nhiên của cơm thôi ạ.''

Hai người cùng gật gù đồng tình, chưa kịp mở lời tranh luận xem ai thắng thì\dots

*GẮP* *NHAI* *GẮP* *NHAI*

Đôi đũa trong tay Suzuno bỗng biến thành một vệt mờ. Với cái bụng ``không đáy'' nổi tiếng của mình, cô bé chẳng hề ngần ngại mà gắp liên tục hết khoanh này đến khoanh khác. Chưa đầy hai phút trôi qua, cả hai đĩa cơm cuộn to tướng đã sạch bách, chẳng còn lại một mẩu nhỏ nào.

``Cảm ơn hai chị vì bữa ăn ngon lắm ạ\~{}'' Suzuno xoa cái bụng vẫn phẳng lì, đẩy hai chiếc đĩa trống sang một bên rồi nở nụ cười thiên thần. Cô quay sang túi rau củ vừa mua, lại còn ngân nga: ``Em chuẩn bị nấu bữa tối cho cả nhà đây\~{}''

Mikeneko và Delutaya đứng chết trân giữa bếp, mắt tròn thao nhìn hai chiếc đĩa trống không. Họ vốn bỏ cả nửa tiếng đồng hồ ra để làm cơm cuộn, định gây ấn tượng với ông anh Cả nhà Hikawa, vậy mà giờ đây chẳng còn gì để trình diện.

Đúng lúc đó, Kuon vừa mở cửa đi vào nhà, cầm trên tay ly nước lọc sau khi tan làm về. Cậu liếc qua hai chiếc đĩa sạch bóng trên bàn, rồi nhìn sang Suzuno đang mãn nguyện vỗ bụng, liền gật đầu đánh giá một cách tự nhiên: ``Ừm, ngon đấy. Hôm nay có cơm cuộn đổi vị là tốt rồi.''

Nói xong, cậu thong thả nhấp ngụm nước rồi quay lưng đi ra phòng khách, chẳng hề hay biết rằng một cuộc chiến sinh tử vừa mới kết thúc ngay trong căn bếp này.

Cả gian bếp rơi vào một khoảng lặng kỳ quái, chẳng ai nói nên lời.

Mikeneko nhìn sang Delutaya. Delutaya cũng nhìn lại Mikeneko. Hai người đứng ngẩn người ra đó, sự bực bội vì bao công sức bị ăn sạch sẽ bỗng chốc hòa tan vào nỗi buồn cười trước cái kết quá bất ngờ. Khoảng cách giữa hai người dường như cũng tan biến theo.

``Này\dots{}'' Mikeneko lầm bầm, khóe môi cô đang giật giật cố nhịn cười. ``Lần sau\dots\ cô nhớ cho thêm chút muối vào đấy nhé.''

Delutaya quay mặt đi, khẽ hừ một tiếng nhưng khóe miệng cũng nhoẻn lên một nụ cười không che giấu được.

``Còn cô thì nhớ là phải cuộn chặt tay vào đi. Lần sau tôi chắc chắn không thua cô đâu.''

\newpage

\stpt{-Histoire 3-}
\stcap{Mỏ hỗn}
\begin{center}
    {\caviar\fontsize{12pt}{12pt}\selectfont Nhân vật: \emp{ryuuhou2} \emp{mike1}}
\end{center}

Đồng hồ treo tường đã chỉ mười hai giờ rưỡi đêm. Cả căn nhà Hikawa như chìm vào giấc ngủ sâu, chỉ còn tiếng tích tắc đều đặn vang vọng từ phòng khách. Kaien, Suzuno và bé Naomi ở tầng ba đã nhắm mắt nghỉ ngơi từ lâu, căn nhà yên ả đến mức có thể nghe thấy tiếng gió thổi bên ngoài cửa sổ.

Thế nhưng, tại góc phòng khách, nơi có hệ thống cách âm bán phần lắp đặt riêng cho góc máy tính, bầu không khí lại hoàn toàn trái ngược. Ánh sáng xanh nhấp nháy từ màn hình lớn chiếu rọi, đánh thức cả sự ồn ào mà đêm khuya không hề có.

Ryuuhou ngồi chồm hổm trước ghế sofa, hai bàn tay siết chặt chiếc tay cầm chơi game đến mức các ngón tay đều bật trắng. Tất cả đèn trong phòng đã bị tắt, chỉ còn ánh sáng màn hình hắt lên gương mặt cô nàng đang bừng bừng lửa giận. Em gái thứ tư của nhà Hikawa đã bước vào chế độ ``quỷ nhỏ'' cực hình, năng lượng tiêu cực trào ra khỏi người chẳng cần ai đốt cháy.

``\textbf{Đèo mẹ, đội này có chơi được không hả?! Đứng đần mặt ra đó làm bia tập bắn thật à?!}'' Ryuuhou gào nhẹ trong kẽ răng, mắt trợn ngược lên màn hình, các ngón tay đập liên tục vào các nút bấm. ``\textbf{Cứ giao tranh mà không có team play, ném hết vào sông cho đỡ tốn diện tích à! Tiến lên bọc lót giùm cái đi lũ vô dụng này!}''

Mikeneko tình cờ đi qua hành lang tầng ba, định xuống bếp lấy mỳ ly ăn đêm. Nghe tiếng bàn phím lạch cạch và tiếng càu nhàu quen thuộc, bản năng của một ``game thủ'' trong cô nàng lập tức thức tỉnh. Nàng Mèo hồng tò mò kiễng chân, rón rén bước vào góc phòng để xem sao.

Đứng sau lưng Ryuuhou được mười giây, chứng kiến cô nàng vừa bị đối phương hạ gục hồi sinh, Mikeneko không kìm được mà cất giọng phàn nàn. ``Em đứng vị trí đó sai bét ra rồi! Phải chiếm cứ điểm giữa map để cắm mắt trước, đâm đầu vào đó khác gì tự sát đưa điểm cho đội bạn đâu!''

Ryuuhou khựng lại, quay ngoắt đầu về phía sau với ánh mắt đỏ ngầu như sém bốc cháy. ``Chị không chơi thì im miệng giùm đi! Đứng đằng sau nói nhiều đúng hơn làm à?''

Nói rồi, cô nàng tóc ngắn đen nhánh vơ lấy chiếc tay cầm thứ hai đang nằm trên bàn, quăng thẳng vào người Mikeneko đang đứng ở cửa. ``Hoặc là cầm lấy cái này vào đây chơi cùng, xem chị nói hay hơn làm được không! Muốn nói nhiều thì có thực lực mà nói, chứ đừng có làm người xem đứng đằng sau chõ mồm!''

Mikeneko bắt lấy chiếc tay cầm giữa không trung, nhếch môi đầy kiêu ngạo. ``Thách chị hả? Để em xem chị làm thế nào cho biết mặt!'' Cô kéo ghế gỗ bên cạnh ngồi xuống, và đó chính là sai lầm lớn nhất của cả buổi tối hôm đó.

Khi hai linh hồn cùng mang bản năng "toxic" khi cầm tay cầm gặp nhau, cả căn phòng khách lập tức biến thành một chiến trường thực thụ. Lửa giận bùng nổ từ cả hai phía, chẳng ai chịu nhường ai một chút nào.

``\textbf{Chị di chuyển kiểu gì thế Mèo hồng vụng về?! Đi chậm như rùa mà cũng đòi carry team hả?!}'' Ryuuhou gào lên, tay bấm liền hồi các nút. ``\textbf{Bọc lót sau lưng cho chị đi, đang bị người ta dính sát rồi!}''

``\textbf{Em mới là đứa đâm đầu vào giữa năm kẻ địch kia đồ ngốc! Muốn chết sớm thì tự đi mà chết, đừng kéo chị theo!}'' Mikeneko đáp trả không kém phần đanh đá, hai bàn tay di chuyển trên tay cầm nhanh đến mức thành vệt mờ. ``\textbf{Dùng chiêu cuối ngay đi, còn chần chờ gì nữa đồ đầu đất!}''

``\textbf{Im mồm đi! Đang hồi chiêu đây! Cứ lèm bà lèm bèm thế này ai mà chơi được game hả?!}''

Thật kỳ lạ thay, dù miệng thì chửi rủa không ngừng, mọi lời lẽ độc hại đều trút vào mặt nhau (và đã được biên tập lại nhẹ nhàng cho phù hợp với cuốn sách này), cách phối hợp của hai người trên màn hình lại ăn khớp đến mức đáng sợ. Mikeneko luôn đứng sau bọc lót cực kỳ chuẩn xác cho những pha lao lên liều lĩnh của Ryuuhou, còn Ryuuhou thì dứt điểm gọn gàng mọi kẻ địch mà Mikeneko đã khống chế. Tiếng cãi vã không làm gì cản trở được sự kết nối ăn ý giữa hai game thủ.

Tiếng la hét và tiếng đập tay cầm ngày một lớn dần, thậm chí còn xuyên qua được những bức tường dày của tầng ba. Không gì có thể ngăn được cơn cuồng nộ của hai cô gái đang chìm trong trận chiến.

*XOẠCH*

Cánh cửa căn hộ tầng trên đột nhiên bật mở ra. Kaigou bước ra cầu thang trong bộ đồ ngủ in hình mèo, bước chân thình thịch bước xuống. Cô đứng trước cửa phòng khách một lúc, hít một hơi thật sâu để bình tĩnh, rồi mới hét lớn vào trong: ``\textbf{Mấy giờ rồi hả?! Hai đứa có chịu im miệng lại để cho mọi người ngủ không?! Có biết mấy giờ đêm rồi không hả?!}''

Nhưng cả Ryuuhou lẫn Mikeneko đều đang chìm đắm trong trận đấu nghẹt thở, tay bấm liên tục nên chẳng hề hay biết gì. Hai người hoàn toàn phớt lờ sự hiện diện đầy đe dọa của cô chị Hai, mắt chỉ dán chặt vào màn hình, miệng vẫn không ngừng chửi rủa những người chơi khác trên mạng.

Kaigou đứng chết trân một phút, nhìn hai đứa trẻ điên cuồng trước mặt mà chỉ có thể thở dài bất lực. Biết mình chẳng thể nói được gì, cô quay lưng bước đi, tay vừa móc điện thoại ra bấm gọi cho anh trai mình. Ngay khi máy vừa nhấc máy, giọng cô than vãn lập tức vang lên. ``Aniki, xuống giúp em với! Hai đứa kia ở tầng ba lại đánh nhau với máy tính rồi, em không dằn được đâu!''

Chưa đầy năm phút sau, Kuon vắt chân chữ ngũ bước xuống cầu thang, tay vẫn cầm tập hồ sơ nhân sự dở dang của công ty. Cậu vừa xoa cằm vừa tiến lại góc phòng máy tính, đúng lúc Ryuuhou vừa ấn nút hồi chiêu cho nhân vật của mình để lao vào giao tranh nữa.

Màn hình đồng hồ đeo tay của Kuon khẽ lóe sáng, giọng nói điện tử từ Game vang lên nhẹ nhàng, đủ để hai cô gái trong phòng nghe thấy. ``Nửa đêm nửa hôm rồi mà hai cô vẫn chưa ngủ à? Định đánh game đến sáng hả? Có để cho người khác ngủ không vậy?''

Cậu anh cả khẽ thở dài, gấp tập hồ sơ lại trên tay. Cậu đứng dậy, bước từ từ lại gần góc máy tính, không hề tỏ ra giận dữ hay la mắng gì cả.

Không cần phải la hét, không cần phải quát mắng. Kuon chỉ cần vươn tay ra, nhẹ nhàng đặt lòng bàn tay rộng lớn của mình lên vai Ryuuhou, hành động đầy dịu dàng và bình thản.

``Ryuuhou này. Ổn không?''

Ngay lập tức, luồng khí nóng hừng hực bốc quanh người cô em gái như tan biến sạch sẽ. Ryuuhou khựng lại, lồng ngực hít vào một hơi thật sâu rồi từ từ thở ra. Ánh mắt dữ tợn ban đầu lập tức dịu lại, cả vai cô thả lỏng hoàn toàn, chẳng còn chút hung hăng nào nữa.

``\dots Em\dots\ em ổn rồi, Onii,'' Ryuuhou nhỏ giọng đáp, trở lại vẻ bình thường một cách kỳ diệu như thể chưa từng có chuyện gì xảy ra.

Mikeneko đứng hình ngay tại chỗ, há hốc miệng nhìn cảnh tượng trước mắt. Cô không thể tin nổi vào mắt mình—một kẻ vừa mới chửi thề như hát, còn đe dọa sẽ đập nát màn hình vài giây trước đó, lại có thể lập tức ngoan ngoãn như cún con chỉ nhờ một cái chạm tay nhẹ nhàng của ông anh trai.

Kuon xoa nhẹ đầu Ryuuhou một cái, rồi quay sang nhìn Mikeneko, giọng nói cũng dịu dàng không kém: ``Muộn rồi đấy. Chơi nốt ván này rồi cả hai nghỉ ngơi đi nhé. Mai Ryuuhou còn phải đến trường học nữa.''

``Vâng ạ\~{}'' cả hai cùng đồng thanh đáp, giọng nói đều nhỏ nhắn như trẻ con vừa bị mắng.

Nói xong, Kuon quay lại bàn thu dọn tập hồ sơ, Kaigou cũng thở phào nhẹ nhõm rồi quay về phòng ngủ của mình. Cô vẫn còn lẩm bẩm không hiểu sao ông anh trai lại có thể dỗ được cô em gái mỏ hỗn kia một cách dễ dàng thế.

Mikeneko lén nhìn sang Ryuuhou, rồi lại nhìn chiếc tay cầm trong tay mình. Bất giác, cô nhận ra một điều\dots\ Cảm giác được xả hết mọi sự bực bội trong ngày, được chửi rủa thoải mái và chơi game hết mình với một người "đồng điệu" như Ryuuhou lại khiến cô thấy thoải mái lạ thường. Không có sự soi xét, không có sự giả tạo hay phải giữ kẽ gì cả, chỉ có hai người cùng thả lỏng bản thân.

Màn hình máy tính bắt đầu đếm ngược, chuẩn bị chuyển sang ván đấu mới. Hai người ngồi sát bên nhau trong im lặng hoàn toàn. Không ai nói thêm câu nào, không còn lời chửi thề hay cãi vã gì nữa. Bầu không khí căng thẳng ban đầu đã tan biến hết, chỉ còn lại sự thấu hiểu ngầm và sự đồng điệu kỳ quặc giữa hai tâm hồn game thủ.

\newpage

\stpt{-Histoire 4-}
\stcap{Hai đứa trẻ}
\begin{center}
    {\caviar\fontsize{12pt}{12pt}\selectfont Nhân vật: \emp{naomi2} \emp{mike1} \emp{koneko}}
\end{center}

Gió chiều se lạnh thổi vi vu qua ban công tầng tám, nơi không gian thoáng đãng nằm ngay sát phòng ngủ mới của Mikeneko, mang theo hương hoa sữa từ hàng cây dưới phố len lỏi vào.

Naomi ngồi một mình trên chiếc ghế gỗ dài ven lan can, đôi chân nhỏ nhắn đung đưa nhịp nhàng theo từng cơn gió thoảng qua. Đôi mắt màu xanh lục bảo trong vắt của cô bé dõi theo những đám mây trắng trôi lững lờ trên bầu trời phố thị, giữ thói quen quan sát vạn vật của một linh hồn đã hơn một nghìn năm tuổi.

Tiếng bước chân nhẹ nhàng như tiếng mèo đi trên thảm len vang lên từ cửa phòng. Koneko vừa hoàn thành việc dọn dẹp góc học tập mới, nên bước ra để hít thở chút không khí trong lành. Nhìn thấy Naomi ngồi đơn độc, đôi tai mèo đen nhánh trên đầu cô bé khẽ rung rung tỏ ý vui mừng. Cô bé lẳng lặng tiến lại, rồi tự nhiên ngồi xuống khoảng trống bên cạnh, chẳng cần phải chào hỏi gì phức tạp.

Cả hai đứa trẻ vốn đều có tính cách hướng nội, chỉ là theo những cách khác nhau. Khoảng im lặng kéo dài giữa hai đứa chẳng hề gượng gạo, mà ngược lại, mang lại cảm giác bình yên khó tả, như thể hai người đã quen nhau từ lâu.

Sau vài phút ngắm nhìn cùng một mây trời, Naomi khẽ xoay người sang, giọng trong trẻo như tiếng nước chảy cất lên: ``Koneko-chan, cậu thấy ở đây thế nào?''

Koneko ngẩn người một chút trước sự chủ động của cô em út nhà Hikawa. Cô bé nhấp môi suy nghĩ mấy giây, đôi mắt hai màu nhìn ra dãy nhà cao tầng mờ ảo trong sương chiều, rồi nhỏ nhẹ nói thật lòng: ``Khác lắm. Nhà của hai chị em mình trước giờ thường rất yên tĩnh\dots\ Ở đây ồn ào hơn nhiều.''

Naomi chỉ gật gật đầu, không phản bác gì. Cô bé nhìn sâu vào mắt Koneko bằng ánh nhìn tĩnh lặng, như thể có thể xuyên thấu được hết những tâm tư chưa nói ra. Nhờ khả năng nhạy cảm tâm linh của linh hồn 1300 năm tuổi, em dễ dàng cảm nhận được những rung động chân thật nhất từ người đối diện.

``Nhưng tớ thấy cậu không ghét sự ồn ào này lắm đâu.''

Koneko há hốc miệng ngạc nhiên, hai tai mèo trên đầu dựng thẳng lên vì bất ngờ. Cô bé chẳng bao giờ ngờ một cô bé nhỏ nhắn trông mới chỉ tám tuổi lại có thể nhìn thấu suy nghĩ của mình rõ ràng đến thế.

Chút ngượng ngùng thoáng qua trên mặt, bé Mèo con cúi đầu, giọng nhỏ nhẹ như muốn nói với chính mình: ``Mình\dots\ mình chỉ là chưa quen thôi.''

Ngay lúc đó, Mikeneko bước ra ban công, định gọi em gái vào uống trà chiều. Nhìn thấy Koneko và Naomi ngồi sát bên nhau, trò chuyện nhỏ, cô nàng bèn dừng bước ở cửa. Không muốn làm phiền khoảnh khắc bình yên của hai đứa trẻ, nàng Mèo hồng chỉ lặng lẽ tiến lại, ngồi xuống chiếc ghế khác phía sau lưng họ, giữ khoảng cách vừa đủ để hai bé thoải mái.

Gió chiều hất tung mấy lọn tóc hồng của nàng Mèo, làm lọn tóc rủ xuống trước trán. Naomi xoay người lại, đôi mắt lục bảo lướt qua Mikeneko một cái rồi lại nhìn Koneko. Đột nhiên, em hỏi một câu nghe có vẻ hồn nhiên, nhưng lại làm không khí xung quanh như dừng lại một nhịp: ``Hai chị em cậu có sợ ở đây không?''

Câu hỏi như chạm thẳng vào những góc khuất nhất trong tâm can hai chị em. Hai chị em nhà Mèo bất giác quay sang nhìn nhau, trong mắt cả hai cùng hiện lên những mảng ký ức cũ: những ngày tháng bị dư luận dị nghị, những đêm không ngủ được vì sợ mai này sẽ phải đi đâu nữa, cảm giác chông chênh như không có chỗ dựa dẫm.

Mikeneko thở dài, thả lỏng toàn bộ bờ vai vốn luôn luôn căng cứng đề phòng. Nụ cười kiêu ngạo, miệng lưỡi đanh đá thường ngày hoàn toàn biến mất, thay vào đó là một vẻ mặt yếu đuối hiếm hoi. Cô ngước nhìn bầu trời hoàng hôn nhuộm đỏ rực, giọng nhẹ hẫng như sợ gió cuốn đi câu nói của mình: ``Không. Hai chị em chị\dots\ chẳng còn chỗ nào khác để mà sợ nữa rồi.''

Naomi im lặng gật đầu, như thể đã hiểu rõ mọi ẩn sâu trong câu nói đó. Cô bé luồn tay vào túi áo khoác bông, rút ra ba cây kẹo mút màu đỏ hồng vị dâu tây, món quà vặt em luôn mang theo bên mình để chia sẻ với những người cô thương.

Em xòe bàn tay nhỏ nhắn, đưa cho Koneko một cây, rồi với tay đặt vào lòng bàn tay Mikeneko một cây khác, giọng nói trong trẻo làm tan đi mọi mây mù: ``Ăn kẹo ngọt sẽ hết sợ đấy. Onii-chan dạy em như thế ạ.''

Mikeneko nhìn cây kẹo trong tay, rồi ngước nhìn gương mặt ngây thơ nhưng ấm áp như mặt trời mùa xuân của Naomi. Đôi mắt cô nàng ươn ướm, nhưng cô vội bóc vỏ kẹo, cho vào miệng để che đi cảm xúc đang trào dâng.

``Ngọt thật đấy\dots'' nàng Mèo khẽ thì thầm, vị ngọt lan tỏa từ đầu lưỡi vào tim, làm tan đi những nỗi sợ cũ.

Cô em gái cũng làm theo, bóc vỏ kẹo rồi ngậm lấy, vị ngọt dâu tây lan tỏa làm cô bé mỉm cười rạng rỡ, câu nói chân thành thoát ra: ``Cảm ơn cậu nhé, Naomi-chan.''

Ánh nắng hoàng hôn dát một lớp vàng ươm lên toàn bộ góc ban công tầng tám. Ba người họ - hai chị em nhà Mèo đã trải qua biết bao sóng gió cuộc đời, và cô bé mang linh hồn cổ xưa hiểu hơn ai hết về nỗi cô đơn - ngồi yên lặng bên nhau. Không ai nói thêm lời nào, chỉ còn tiếng gió thổi qua lan can, vị ngọt của kẹo còn đọng lại trong miệng, cùng nhau ngắm nhìn thành phố dần lấp đầy ánh đèn lung linh khi màn đêm buông xuống.

\newpage

\stpt{-Histoire 5-}
\stcap{Sự ngăn nắp}
\begin{center}
    {\caviar\fontsize{12pt}{12pt}\selectfont Nhân vật: \emp{suzuno2} \emp{koneko}}
\end{center}

Ánh đèn vàng ấm áp len lỏi khắp từng góc căn bếp tầng sáu, bao trùm mọi thứ sau khi bữa tối muộn vừa khép lại. Tiếng nước róc rách chảy từ vòi rửa hòa quyện cùng tiếng chén đĩa va chạm nhẹ nhàng, tạo nên một bản giao hưởng bình dị của gia đình.

Gió từ cửa sổ nhỏ mang theo hương trà thảo mộc đang ấm trên bàn, làm xao động vài lọn tóc thả lỏng của người đứng bên bồn rửa. Suzuno đang cẩn thận tráng từng chiếc bát, đôi tay đeo trong đôi găng tay cao su màu hồng xinh xắn. Mỗi khi em cử động, ba chiếc chuông vàng nho nhỏ gắn trên chiếc vòng tay đỏ kỷ vật lại khẽ leng keng, vang lên những âm thanh êm dịu như đang hát ru cho căn bếp.

Tiếng bước chân nhỏ nhẹ vang lên từ hành lang, Koneko bước vào bếp với đôi tai mèo đen nhánh khẽ rung rung. Cô bé đã quen với việc luôn giữ gìn sự ngăn nắp cùng chị gái Mikeneko, nên không cần ai nhờ vả, tự nhiên tiến về phía kệ chén. Bé Mèo con lặng lẽ cầm lấy chiếc khăn bông mềm khô, đỡ từng chiếc đĩa vừa được tráng sạch để lau khô rồi xếp chúng lên chạn.

Không hề có lời giao tranh nào vang lên, cả hai chỉ phối hợp làm việc song song trong sự im lặng nhẹ nhàng như đã quen thuộc từ lâu. Chỉ có tiếng nước chảy và tiếng khăn giấy nhẹ nhàng lau chùi, xen lẫn tiếng chuông trên tay Suzuno thỉnh thoảng leng keng.

Koneko cẩn thận phân loại bát đĩa theo một logic rất riêng của mình: tất cả những chiếc bát màu hồng được xếp chung một chồng, bát màu đen gom vào một chồng khác, còn những chiếc bát sứ trắng tinh thì để riêng ở một góc. Cô bé cứ xếp vừa nhìn vừa mỉm cười nhỏ, hài lòng với sự đồng bộ màu sắc sắp tạo ra.

Đang tráng nốt chiếc chảo lớn cuối cùng, Suzuno vô tình liếc qua cách Koneko sắp xếp và khẽ bật cười. Em tháo đôi găng tay cao su ra, lau khô tay rồi tiến lại gần kệ chén. Không hề làm đổ vỡ gì, em chỉ nhẹ nhàng nhấc toàn bộ chồng bát đĩa Koneko vừa xếp ra, rồi bình tĩnh sắp xếp lại: đĩa lòng sâu đặt xuống dưới cùng, đĩa phẳng xếp lên trên, còn bát ăn cơm thì được phân loại theo kích cỡ to nhỏ thay vì màu sắc.

Koneko chớp chớp đôi mắt hai màu to tròn, hai tai mèo trên đầu khẽ vẫy ngược lại vì bất ngờ. Cô bé ngước nhìn cô chị Ba với ánh mắt tò mò không che giấu, rồi nhỏ giọng hỏi: ``Sao chị lại xếp theo kích cỡ ạ? Xếp theo màu sắc không phải trông sẽ đẹp mắt hơn sao?''

Suzuno quay sang nhìn cô bế, đôi mắt dị sắc một đỏ một xanh phản chiếu ánh đèn bếp vàng ấm, trông thật dịu dàng. Em không hề có vẻ gì là đang sửa lỗi hay dạy dỗ, mà chỉ mỉm cười giải thích bằng giọng nói ngọt ngào, gần gũi như một người chị ruột.

``Nếu xếp theo màu thì nhìn rất thích mắt, nhưng lúc lấy ra dùng lại hơi bất tiện đấy Koneko-chan. Bếp nhà mình đông người, mỗi người có thói quen ăn uống khác nhau. Nếu đặt đĩa sâu lòng ở nơi dễ lấy nhất, mọi người sẽ không phải loay hoay tìm đồ đựng thức ăn có xốt.''

Koneko chú ý lắng nghe từng lời, đầu khẽ gật gù như đang ghi nhớ vào lòng. Cô bé lướt mắt quanh toàn bộ gian bếp một lượt, mới sực nhận ra mọi thứ ở đây đều có vị trí của nó: từng chiếc hũ gia vị, từng chiếc đĩa, khay gỗ hay cả đũa muỗng đều được đặt ở những vị trí cực kỳ chính xác. Không có bất kỳ sự ngẫu nhiên nào ở đây cả, mọi thứ đều được sắp xếp dựa trên sự thấu hiểu thói quen của từng thành viên.

Cô bé đứng im một lát, rồi mới nhỏ giọng hỏi với chút ngưỡng mộ lẫn ngạc nhiên: ``Suzuno-oneechan sắp xếp mọi thứ theo thói quen của từng người trong nhà sao ạ?''

Suzuno gật đầu ngay, ánh mắt ánh lên sự săn sóc tự nhiên như một người mẹ chăm lo cho cả nhà. Em lại gần Koneko, giọng nói vẫn ấm áp như mọi khi khi liệt kê từng người.

``Ừm! Onee-sama hay nấu mấy món hầm lớn nên nồi lớn luôn được để ở ngăn dưới cùng dễ lấy nhất. Onii-sama thì buổi sáng hay vội đi làm, nên khay ăn nhanh và ly cà phê của anh ấy luôn ở ngay tầm tay. Còn Ryuuhou-chan thì\dots'' cô khẽ bật cười hờn dỗi đáng yêu, vai khẽ rung lên vì nhớ ra chuyện buồn cười. ``\dots em ấy hay ôm bát đĩa đi chơi game rồi bỏ quên khắp nơi, nên chị phải cất bớt một nửa đũa riêng của cậu ấy vào tủ, khi nào gom đủ mới đưa ra lại.''

Koneko im lặng lắng nghe, lòng dâng lên một cảm xúc ấm áp không thể nói nên lời. Sự tỉ mỉ và tình yêu thương thầm lặng mà Suzuno dành cho căn nhà này khiến trái tim cô bé xao xuyến. Koneko nhìn xuống đôi bàn tay nhỏ nhắn của mình, lầm bầm vài giây như đang tập hợp sự can đảm, rồi mới ngước lên hỏi bằng một giọng rất khẽ, chứa đựng chút ngập ngừng sợ hãi: ``Vậy\dots\ chị có thể sắp xếp thêm chỗ cho Onee-chan và em được không ạ?''

Suzuno không cần mất một giây suy nghĩ nào cả, em chỉ mỉm cười dịu dàng và bước sang bên cạnh Koneko. Em kéo nhẹ một ngăn kéo gỗ phụ ngay dưới mặt bếp, mở rộng ra để cô bé có thể nhìn rõ mọi thứ bên trong.

Trong ngăn kéo sạch sẽ khô ráo ấy, hai bộ bát đĩa sứ mới tinh—một bộ viền hồng đào xinh xắn và một bộ viền đen tuyền tối giản—đã được xếp ngăn nắp từ bao giờ. Ngay bên cạnh là hai đôi đũa gỗ được làm thủ công, được bọc trong giấy gói nhỏ còn nguyên vẹn. Nhìn chúng thôi cũng biết đã được chuẩn bị rất chu đáo.

``Chị đã chuẩn bị sẵn chỗ cho hai chị em từ trước khi em hỏi rồi đây này,'' Suzuno nghiêng đầu mỉm cười thiên thần, ánh mắt tròn đầy yêu thương khi nhìn cô bé.

Koneko đứng yên nhìn vào ngăn kéo gỗ, đôi mắt hai màu của cô bé hơi ươn ướm. Những chiếc bát đĩa nằm đó như một lời khẳng định thầm lặng rằng hai chị em cô không phải là những kẻ ở nhờ tạm bợ, mà đã thực sự trở thành một phần của mái nhà ấm áp này. Mỗi chi tiết nhỏ đều dành riêng cho họ, không phải là đồ đạc tạm thời để dùng qua loa.

Cô bé nhìn ngăn kê một lúc lâu, rồi ngước lên nhìn nụ cười ấm áp của Suzuno. Koneko không thốt thêm lời cảm ơn màu mè nào, vì cô biết những lời đó chẳng đủ diễn tả lòng mình. Cô bé chỉ khẽ gật đầu đầy kiên định và thở phào nhẹ nhõm của một người vừa nhận ra mình đã tìm đúng nơi thuộc về.

\newpage

\stpt{-Histoire 6-}
\stcap{Nơi chào đón}
\begin{center}
    {\caviar\fontsize{12pt}{12pt}\selectfont Nhân vật: \emp{kuon2} \emp{mike1} \emp{koneko}}
\end{center}

Đêm đã chuyển sang khuya lắm rồi. Cả căn nhà Hikawa chìm vào một sự yên tĩnh đến lạ thường, chỉ còn ánh đèn ngủ vàng ấm hắt hơi qua khe cửa hành lang tầng tám, vẽ những vệt sáng mỏng manh trên sàn gỗ.

*CỐC CỐC*

Tiếng gõ cửa nhẹ nhàng vang lên đúng hai nhịp, đủ để phá tan sự im lặng mà không làm phiền ai. Kuon đẩy cửa một chút rồi bước vào, nhằm chỉ để lướt qua kiểm tra xem cả hai đã quen với phòng mới chưa.

Trong phòng, Mikeneko trông đã tự nhiên hơn hẳn so với những ngày đầu. Cô nàng đang nằm dài bò trên giường, hai chân trắng muốt đung đưa qua lại, tai đeo cặp tai nghe trắng thưởng thức những giai điệu rock ầm ĩ mà cậu chỉ nghe được tiếng rì rào nhỏ qua lớp vỏ.

Ở góc phòng, chiếc bàn học nhỏ được bày biện gọn gàng, Koneko đang ngồi quay lưng lại với cửa. Đôi tai mèo đen nhánh trên đầu cô bé khẽ rung lên khi nghe tiếng bước chân quen thuộc. Cô bé từ từ đặt nắp bút máy lên, gấp gọn tờ giấy đang viết rồi ngẩng lên nhìn Kuon, đôi mắt hai màu bình thản như đã chờ đợi từ lâu.

``Mọi thứ trong phòng vẫn ổn chứ?'' cậu Tổng hỏi, giọng nói trầm ấm như mọi khi, không hề lớn tiếng để làm phiền không gian yên tĩnh.

Nàng Mèo hồng vội tháo tai nghe ra, ngồi dậy nhanh chóng và gật đầu lia lịa: ``Ổn cả ạ! Phòng này rộng rãi hơn chỗ cũ của bọn em nhiều lắm.''

Thế nhưng Koneko lại không theo lời chị mình. Cô bé xoay ghế lại đối mặt với Kuon, đặt hai tay lên đỉnh đầu gối rồi cất lên câu hỏi đã ấp ủ trong lòng từ ngày đầu tiên đặt chân vào căn nhà này: ``Onii-chan\dots\ tại sao chị em chúng em lại phải chui vào cái thùng các-tông lớn đó để anh khiêng lên, thay vì bấm chuông đứng ở cửa chờ như mọi người bình thường ạ?''

Không khí trong phòng lập tức đọng lại như nước trong bình vừa ngừng chuyển động. Cậu Tổng khựng lại một nhịp, khóe mắt thoáng lên một chút ngạc nhiên rồi cậu khẽ vuốt cằm: ``\dots Anh cũng đang định hỏi như thế. Tại sao?''

Mikeneko nghe vậy thì giật nẩy người, vội bật dậy ngồi thẳng lưng, hai bàn tay siết chặt vào mép giường. Cô nàng vội vàng tuôn ra một tràng lý lẽ đã chuẩn bị từ trước, giọng nói hơi cao lên để che đi sự bối rối đang lan tỏa trên mặt: ``Bởi vì\dots! Em đã tính toán kỹ lắm! Nếu đứng trước cửa nhà anh mà bấm chuông thì nguy cơ bị hàng xóm hay người qua đường nhận ra là rất cao! Với lại\dots\ cái hộp các-tông đó còn bảo vệ bọn em khỏi ánh mắt soi mói của mọi người nữa mà--''

``Onee-chan,'' cô em gái điềm tĩnh cắt ngang lời chị, ánh mắt không hề có chút dao động nào. ``Giải thích như thế anh ấy sẽ không tin đâu. Chị thừa biết con phố này muộn thế nào cũng chẳng có một bóng người qua lại mà.''

Cô chị lập tức cứng họng, những lời nói còn đang chảy trong miệng bỗng chốc tắc nghẹn lại nơi cổ, không thể thốt ra thêm một chữ nào.

Kuon nhìn cô nàng Mèo hồng đang bối rối, cố gắng giữ cho nét mặt mình thật thản nhiên để không làm cô thêm áp lực. Cậu kéo chiếc ghế gỗ đối diện giường ngồi xuống, khoanh tay lại với ánh mắt bao dung nhưng như muốn xuyên thấu vào sự thật đang bị che giấu: ``Vậy\dots\ lý do chính là gì thế, Mikeneko?''

Sự im lặng kéo dài lâu hơn bao giờ hết. Mikeneko cúi gằm mặt xuống, hai bàn tay siết chặt lấy viền chăn đến trắng bệch. Lớp vỏ kiêu ngạo, miệng lưỡi đanh đá mà cô nàng thường mặc lên người mỗi ngày bỗng chốc bong tróc từng mảng, để lộ ra sự yếu đuối và nỗi sợ hãi đã chôn sâu trong lòng từ rất lâu.

``\dots Vì trước khi chui vào hộp,'' cô nàng nhỏ giọng nói, tiếng nói lí nhí nhỏ đến mức gần như biến mất vào không khí, ``em đã đứng trước cửa nhà anh khoảng\dots\ hai mươi phút. Em không dám bấm chuông. Em cứ đứng nhìn cái nút chuông đó mãi mà tay chẳng thể nhấn xuống được.''

Cô nàng ngước mắt lên, viền mắt đã khẽ đỏ lên vì xúc động.

``Nếu em bấm chuông và đứng đó ở cửa\dots\ em sẽ phải đối mặt. Và nếu anh không xuống, em cũng chẳng biết mình phải đi đâu nữa. Ít nhất\dots\ nếu chui vào trong hộp, thì không ai có thể từ chối một món hàng đã được giao đến tận nơi. Em\dots\ em sợ phải là người đầu tiên cất lời hỏi xin ở lại.''

Mikeneko ngập ngừng thêm một chút rồi nói tiếp, một nụ cười cay đắng hiện lên nơi khóe môi: ``Lúc trước\dots\ Coco-chan\dots\ à không, Kson-chan\dots\ cái đêm em ấy chưa trở thành Kson, em ấy đã gặp em trên vách đá\dots. Lúc đó, emấy từng hỏi em một câu: `Chị nghiêm túc đấy chứ? Chị thực sự đành lòng để họ khóc vì chị như thế kia sao?' Em đã không thể trả lời được. Và từ đó\dots\ em sợ phải đối mặt với bất kỳ câu hỏi nào tương tự.''

Koneko ngồi bên cạnh nhìn chị mình, sau đó nhẹ nhàng nói thêm một câu mà chỉ có cô mới hiểu rõ sự thật nhất: ``Onee-chan từng bảo em: `Nếu anh ấy không mở hộp ra thì chúng ta sẽ tự mở rồi đi về thôi.' Nhưng em biết chị ấy chưa bao giờ tin vào câu đó cả. Chị chỉ muốn có một người bắt buộc phải mở chiếc hộp ra đón nhận mình, để chị không cần phải tự mình cất lời xin xỏ bất cứ điều gì.''

Cả căn phòng lại chìm vào tĩnh lặng hoàn toàn. Màn hình đồng hồ thông minh trên tay Kuon bỗng lóe sáng một chút, nhưng trợ lý Game lần này lại tuyệt đối im lặng, không hề đưa ra bất kỳ bình luận hay phân tích nào. Một sự im lặng tôn trọng hiếm thấy từ vị trợ lý ảo kỹ thuật số.

Kuon nhìn nàng Mèo hồng đang cúi đầu, bờ vai nhỏ của cô nàng khẽ run lên vì đang cố gắng kìm nén cảm xúc. Cậu không nói những lời triết lý suông, cũng chẳng la mắng cô nàng vì sự ngây thơ của nỗi sợ. Cậu chỉ đứng dậy, bước lại gần giường rồi vươn tay vỗ nhẹ lên mái tóc hồng của cô, đúng như cách cậu vẫn thường làm để vỗ về bé Naomi mỗi khi em ấy gặp ác mộng. 

``Lần sau cứ bấm chuông. Anh sẽ mở cửa.''

Câu nói ngắn gọn, không hề có bất kỳ điều kiện nào đi kèm, đã trả lời trực tiếp vào nỗi sợ lớn nhất trong lòng cô nàng. Mikeneko không cần phải dùng bất kỳ chiêu trò nào để ẩn nấp nữa, vì cánh cửa của ngôi nhà này sẽ luôn luôn rộng mở để đón hai chị em cô.

Cô nàng ngẩn người ngước lên, đôi tai mèo trên đầu khẽ giật giật vì bất ngờ, một luồng ấm áp kỳ lạ dâng tràn lên trong lồng ngực. Cô nàng chỉ khẽ gật đầu, rồi vội cúi mặt xuống để giấu đi giọt nước mắt vừa lăn xuống trên má.

Koneko quan sát trọn vẹn cảnh tượng trước mắt. Cô bé gật đầu nhỏ một mình, sau đó lẳng lặng mở cuốn sổ tay nhỏ trên bàn ra, cầm bút máy viết thêm một dòng chữ ngắn gọn vào trang giấy trắng.

Trước khi bước ra khỏi phòng và khép hờ cánh cửa, Kuon ngoái đầu lại nhìn cô em gái nhỏ đang cầm bút: ``Em đang ghi gì vào sổ thế, Koneko?''

Bé Mèo con nhẹ nhàng đóng cuốn sổ lại, ngước nhìn cậu với một nụ cười hiếm hoi thoáng qua trên môi: ``Dạ, em chỉ ghi những thứ em cần nhớ về nơi này thôi ạ.''

Kuon mỉm cười rồi rời đi, để lại không gian yên bình ấm áp cho hai chị em trong phòng.

Trên trang sổ vừa mở của Koneko, bên dưới những dòng ghi chú tỉ mỉ về giờ giấc sinh hoạt của từng người trong nhà Hikawa, là một dòng chữ mực đen còn mới nguyên, được viết bằng tất cả sự tin tưởng và bình yên mà cô bé chưa từng cảm thấy trước đây:

\begin{center}
    ``Ở đây: Cánh cửa luôn chào đón chúng ta.''
\end{center}

\end{document}